% Versión en español, autocontenida (conceptos explicados al usarse) y acotada.
% Compilar con pdfLaTeX (Overleaf) o tectonic. Las figuras viven en figures/.
\documentclass[11pt]{article}
\usepackage[a4paper,margin=2.3cm]{geometry}
\usepackage{iftex}
\ifPDFTeX
  \usepackage[utf8]{inputenc}
  \usepackage[T1]{fontenc}
\else
  \usepackage{fontspec}
\fi
\usepackage[spanish,es-tabla,shorthands=off]{babel}
\usepackage{graphicx}
\usepackage{amsmath}
\usepackage{booktabs}
\usepackage{caption}
\usepackage[hidelinks]{hyperref}
\graphicspath{{figures/}}
\captionsetup{font=small}

\title{\vspace{-1.2cm}Un retardo distribuido de acumulación en el ciclo ganancias--inversión:\\
identificación de un feedback débil y sus límites}
\author{Juan I. Tollo}
\date{\today}

\begin{document}
\maketitle
\vspace{-0.7cm}

\begin{abstract}
\noindent
Estudiamos el ciclo ganancias--inversión de la economía de EE.UU.\ con datos trimestrales
(1948--2026). Partimos de la hipótesis de \emph{sobreacumulación}: la inversión pasada deprime la
rentabilidad presente, pero con un retardo. En los datos encontramos dos cosas a la vez. Por un
lado, las ganancias y la inversión se mueven juntas en el mismo trimestre; ese co-movimiento
contemporáneo es el rasgo dominante. Por otro, aparece un feedback negativo más débil, concentrado
en torno a un año. Esta combinación descarta a los osciladores depredador--presa de \emph{fase
fija}, como Lotka--Volterra o FitzHugh--Nagumo. En ellos el desfase entre las dos variables es fijo
---un cuarto de ciclo, por construcción---, lo que fuerza una correlación contemporánea \emph{nula},
$\rho(0)=0$, justo lo contrario de la fuerte correlación contemporánea que observamos,
$\rho(0)\approx+0{,}5$. Proponemos en cambio un modelo físico de \emph{cadena de acumulación}: una
ecuación diferencial con retardo distribuido. Identificamos el \emph{centro} de ese retardo ---$\sim$5 trimestres, alrededor de un año--- con un
inverso bayesiano que lo estima con su incertidumbre, cotejado con la correlación cruzada; su
\emph{distribución}, en cambio, no se identifica. El mismo posterior concentra su
masa en un coeficiente de sobreacumulación negativo ($k<0$): la inversión madurada deprime la
ganancia posterior. El efecto y su retardo de $\sim$1 año \emph{sí} se identifican y son robustos
---no se deben a las crisis extremas (2008/2020) ni a la transformación de la serie---.
El resultado central es de \emph{identificabilidad}: con una señal que explica apenas un 3\% de la
varianza, el centro del retardo y el \emph{signo} del efecto se fijan, pero no su distribución ni su separación de un
$\mathrm{VAR}(p\!\ge\!4)$. Que el efecto sea débil en varianza no lo vuelve irrelevante: su dirección ---una presión de sobreacumulación sobre la ganancia--- es consistente y estructural.
\end{abstract}

\section{Introducción}

\paragraph{La hipótesis económica.} Las ganancias y la inversión forman un ciclo endógeno. Las
ganancias altas estimulan la inversión, pero esa misma inversión, al acumularse, erosiona las
ganancias posteriores. Este vaivén se ve directamente en las series trimestrales de EE.UU.\ (Fig.~\ref{fig:series}). Tapia (2023, pp.~183--186) documenta esta estructura bidireccional para
EE.UU.\ y la formaliza como ``un modelo de tipo depredador--presa'', con las ganancias como presa y
la inversión como depredador: los movimientos de las ganancias son seguidos ``algunos trimestres
después por movimientos de la inversión en la misma dirección, y los movimientos de la inversión son
seguidos por movimientos de las ganancias en la dirección opuesta''. Advierte, además, que esa
segunda dirección (inversión\,$\to$\,ganancias) ---un ``efecto rezagado negativo'' de la variación de
la inversión sobre la de las ganancias--- es la ``más débil'' (traducción nuestra). Astarita (2012) propone leer ese
feedback negativo como \emph{sobreacumulación}: la inversión que, al acumularse en capacidad
instalada, presiona a la baja la rentabilidad ---``una crisis por sobreacumulación, o
sobreinversión'', no por subacumulación---. El antecedente formal de estos osciladores es la
dinámica de Lotka--Volterra (LV), usada en macrodinámica y estimada por Goldstein (1999) para el ciclo del \emph{profit squeeze},
sobre otras variables (no el par ganancias--inversión que estudiamos aquí). En un sistema LV las dos variables oscilan con un \emph{desfase fijo}: el depredador
sigue a la presa con un cuarto de ciclo de atraso, siempre, por construcción.

\paragraph{Nuestra pregunta.} Planteamos una pregunta acotada y mecanística. ¿El feedback de la
inversión hacia las ganancias es un antagonismo \emph{instantáneo}, como imponen esos osciladores?
¿O es un antagonismo \emph{retardado}, cuyo tiempo característico refleja la \emph{acumulación} de la
inversión pasada en capacidad instalada? Encontramos lo segundo. La contribución es descriptiva y
metodológica: una estimación del retardo de acumulación con un inverso bayesiano ---con su
incertidumbre---, cotejada con la correlación cruzada (sección~\ref{sec:metodos}), junto con una
caracterización explícita de lo que estos datos \emph{no} permiten identificar: la forma del retardo
(distribuido vs puntual) y su separación de un VAR de memoria larga. El aporte metodológico, en una
línea: bajo señal débil, recuperar un parámetro estructural es una cuestión de \emph{cómo} se estima,
no solo de cuántos datos hay ---una calibración probabilística directa fija el centro del retardo
donde una red sustituta flexible (PINN) lo disuelve y un VAR en forma reducida no lo separa de sus
rezagos.

\section{Datos y métodos}\label{sec:metodos}

\paragraph{Datos.} Usamos series trimestrales de EE.UU.: ganancias corporativas e inversión privada
bruta (FRED, 1948--2026, $n=313$ trimestres; 312 al pasar a tasas de crecimiento). Trabajamos con tasas de crecimiento \emph{log trimestre a
trimestre} ($x_t=\log Y_t-\log Y_{t-1}$): es con esta diferencia que el lag de $\sim$1 año aparece de
forma estable. La transformación interanual habitual (``YoY'', el cambio respecto del mismo trimestre
del año anterior) no sirve: al ser una diferencia a 4 trimestres, mezcla un año de datos en cada
punto ---la misma escala del feedback que medimos---, así que la ventana del filtro se confunde con
la señal y el lag estimado se corre a $\sim$2 trimestres.

\begin{figure}[t]
\centering
\includegraphics[width=\linewidth]{p15_crisis_series.pdf}
\caption{El ciclo, a la vista. Crecimiento log-trimestral de las ganancias y de la inversión en
EE.UU.\ (en desvíos estándar); las franjas grises son las recesiones del NBER. Las dos series
\emph{co-mueven fuerte y a la vez} ($\approx+0{,}5$) y se desploman juntas en cada recesión; la
ganancia adelanta a la inversión por cerca de un trimestre (la correlación con la inversión del
trimestre siguiente es $+0{,}33$), así que en las caídas la inversión sigue alta cuando la ganancia
ya giró. Todo esto es el \emph{acelerador} (ganancia$\to$inversión, rápido). El feedback que estudia
este trabajo es el \emph{opuesto y más lento} ---la inversión deprime la ganancia $\sim$1 año después
(sobreacumulación)---, más débil y no visible a ojo aquí: se mide con la CCF (valle en el rezago
$-4$) y el inverso bayesiano más abajo.}
\label{fig:series}
\end{figure}

\paragraph{La herramienta descriptiva.} La \emph{función de correlación cruzada} (CCF), definida
como $\rho(h)=\mathrm{corr}(\text{gan.}_t,\text{inv.}_{t+h})$, decide si hay un \emph{ciclo} o solo
co-movimiento. Un valor negativo en $h=-4$ significa que la inversión de hace 4 trimestres se asocia
con ganancias \emph{menores} hoy. Si la CCF oscila ---cruza a negativo en los rezagos--- hay ciclo;
si decae a cero, es co-movimiento puro. Como respaldo, un \emph{estudio de evento} ---alinear las
doce recesiones del NBER (\emph{National Bureau of Economic Research}, que data oficialmente el ciclo
de EE.UU.) y promediar--- corrobora el patrón típico sin que ningún episodio domine.

\paragraph{El modelo físico.} El acelerador ---las ganancias estimulan la inversión--- es
instantáneo y no problemático; lo difícil es la otra dirección, la presión que la inversión ejerce
\emph{con retardo} sobre las ganancias. Ese retardo no es un número fijo: la inversión se construye,
entra en operación y agrega capacidad; cuando esa capacidad produce más de lo que el mercado absorbe
---mercancía que no se realiza---, recién entonces presiona la rentabilidad, de modo que sus efectos
llegan \emph{repartidos} en el tiempo. Lo modelamos como una \emph{cadena} de $N$ etapas
iguales, donde cada etapa $j$ persigue a la anterior con una constante de tiempo $\theta$,
\[
\frac{\mathrm{d}m_j}{\mathrm{d}t}=\frac{m_{j-1}-m_j}{\theta},\qquad m_0=I.
\]
Cada etapa es un suavizado exponencial: su salida $m_j$ es una versión demorada y alisada de su
entrada $m_{j-1}$. La inversión entra por $m_0=I$ y la salida de la última etapa es la
\emph{inversión acumulada} $m$.

Para ver qué hace la cadena, conviene seguir un único pulso de inversión. Tras una sola etapa su
efecto decae como una exponencial $e^{-t/\theta}$; al encadenar $N$ etapas idénticas ---es decir, al
convolucionar $N$ exponenciales--- el efecto deja de estar pegado al origen y pasa a \emph{subir,
hacer pico y bajar}. El ``\emph{truco de la cadena lineal}'', un resultado estándar de las ecuaciones
con retardo (Smith, 2011), dice que esa respuesta es \emph{exactamente} un núcleo Gamma (Erlang) de forma $N$ y
escala $\theta$,
\[
w(t)\;\propto\;t^{\,N-1}\,e^{-t/\theta}.
\]
La lectura probabilística lo vuelve transparente: el tiempo que una unidad de inversión tarda en
atravesar cada etapa es exponencial de media $\theta$, y el tiempo \emph{total} hasta salir de la
cadena es la suma de $N$ de esas demoras independientes. Como en toda suma de variables
independientes, las medias y las varianzas se suman; de ahí que el retardo típico sea $\mu=N\theta$ y
su dispersión $\sqrt{N}\,\theta$. Muchas etapas ($N$ grande) concentran el núcleo alrededor de
$\mu$; pocas etapas lo dejan ancho y asimétrico.

El sistema agregado es entonces
\begin{equation}
\begin{aligned}
\frac{\mathrm{d}P}{\mathrm{d}t} &= c\,I - d\,P + k\,m,\\
\frac{\mathrm{d}I}{\mathrm{d}t} &= a\,P - b\,I,\\
m(t) &= \sum_j w_j(\mu)\,I(t-j),
\end{aligned}
\end{equation}
donde $P$ son las ganancias, $I$ la inversión y $m$ la \emph{inversión acumulada}: un promedio
ponderado de la inversión pasada cuyos pesos $w_j(\mu)$ son el núcleo Gamma, centrados en los
$\sim\mu$ trimestres previos y normalizados a uno. Cada $w_j(\mu)$ es la fracción de la inversión de
hace $j$ trimestres que \emph{recién ahora} termina de pesar sobre la rentabilidad. El parámetro
$\mu$ es el \emph{retardo de acumulación}, lo que queremos estimar; el término $+k\,m$ (con $k<0$) es la
sobreacumulación: la inversión que se acumuló hace $\sim\mu$ trimestres deprime las ganancias de hoy
con intensidad $|k|$. A diferencia del oscilador de Lotka--Volterra clásico ---no lineal, donde el
ciclo nace del \emph{producto} de las dos variables y el desfase queda fijado por construcción---,
este sistema es \emph{lineal} y su oscilación proviene del \emph{retardo}: es el núcleo Gamma, y no
una no linealidad, lo que genera el ciclo, y el desfase lo fijan los datos en lugar de la forma del
modelo.

\paragraph{Estimación del retardo.} Estimamos el retardo $\mu$ con un \emph{inverso bayesiano} del modelo físico ---inferir los
parámetros ($\mu$, $k$) a partir de los datos--- (Turing.jl, NUTS), siguiendo el tratamiento de
problemas inversos de la materia (Sapienza, 2026): en vez de un valor puntual
entrega un \emph{posterior}, una distribución con intervalos creíbles, sobre $\mu$ y sobre el
coeficiente $k$. Lo cotejamos con la \emph{correlación cruzada} (CCF), un método que no supone el
modelo: su valle ubica el retardo de forma independiente. Un inverso \emph{diferenciable} con red
sustituta (PINN) no mejora la identificación; lo discutimos en el Apéndice.

\paragraph{Controles de identificabilidad.} \emph{Identificable} significa que los datos bastan para
fijar un parámetro; si distintos valores ajustan igual de bien, no lo está. Lo evaluamos con cuatro
controles. (i) Un \emph{perfil de verosimilitud} en $\mu$: lo barremos fijándolo y reajustando el
resto, de modo que una curva de ajuste plana indica que $\mu$ no se identifica. (ii) Un \emph{placebo}: repetimos el ejercicio en ventanas sin crisis; si el valle
aparece igual, no es firma de las crisis. (iii) Un \emph{null} por series sustitutas (\emph{surrogates}) AR: simulamos $\sim$2000 pares de
series artificiales que conservan la autocorrelación de cada serie y el co-movimiento contemporáneo,
pero rompen el feedback cruzado \emph{retardado} ---cada serie evoluciona solo con su propio
pasado---. Si el valle observado cae fuera de lo que produce esa nula, no es un artefacto de la
autocorrelación. (iv) La
CCF que \emph{implicaría} un $\mathrm{VAR}(p)$ ---un modelo lineal donde cada variable depende solo
de sus últimos $p$ rezagos---: sirve para ver si un modelo \emph{sin mecanismo} ---sin la cadena de
maduración que postulamos, solo rezagos estadísticos--- ya reproduce el patrón observado.

\section{Resultados}

El patrón es un feedback \emph{retardado}, no un antagonismo instantáneo. Para separar las dos
direcciones del ciclo, la CCF \emph{oscila} en lugar de decaer (Fig.~\ref{fig:ccf}).
La primera dirección es el pico contemporáneo ($\rho(0)\approx+0{,}5$): las ganancias preceden a la
inversión por $\sim$1 trimestre, la dirección fuerte (el acelerador), y que la inversión siga a la
ganancia ---y no la anteceda--- se ve en la asimetría, $+0{,}33$ con la inversión del trimestre
siguiente contra $+0{,}07$ con la del previo, también en las caídas. La segunda es el valle negativo
en el rezago~$-4$ ($\approx-0{,}2$): la inversión de hace $\sim$1 año deprime las ganancias, la
dirección débil de sobreacumulación que estudiamos. Esa forma ---co-movimiento contemporáneo más un
valle negativo retardado---
descarta los osciladores de fase fija (LV, FN; Strogatz, 2015): su desfase fijo de un cuarto de ciclo forzaría
$\rho(0)=0$ con el máximo lejos del origen, lo contrario de lo que muestran los datos.

\begin{figure}[t]
\centering
\includegraphics[width=\linewidth]{p7b_ccf_surrogate.pdf}
\caption{Correlación cruzada ganancias$\leftrightarrow$inversión (log-trim, sin suavizado). La
función oscila: fuerte co-movimiento contemporáneo (rezago 0) y un feedback negativo retardado
a $\sim$1 año (rezago $-4$). La banda gris es un \emph{null} por surrogates AR que preserva la
autocorrelación de cada serie y el co-movimiento contemporáneo, pero destruye el feedback cruzado.
El valle a 4 trimestres cae \emph{fuera} de esa banda ($p\approx0{,}006$, dos colas), de modo que el feedback
es real, no un artefacto de autocorrelación. Es robusto además a excluir 2008/2020.}
\label{fig:ccf}
\end{figure}

El retardo se identifica en torno a un año: el inverso bayesiano lo concentra en $\mu\approx4{,}9$
trimestres $[3{,}0;7{,}3]$, mucho más angosto que el prior, de modo que el dato \emph{informa} $\mu$. Es consistente con el valle de la CCF (en torno a
los rezagos $-4$/$-5$) y con el ajuste del modelo físico de maduración, cuyo $R^2$ es máximo en
$\mu\approx5$ trimestres (Fig.~\ref{fig:model}); $\mu$ es la \emph{media} del núcleo ---integra todo
ese pozo de la CCF, no un rezago aislado---, y por eso cae un poco más allá del rezago más profundo. La identificación emerge al \emph{agrupar} las crisis: por episodio
individual el retardo es demasiado ruidoso para anclarse, pero su centro común se concentra al
combinarlas.

\begin{figure}[t]
\centering
\includegraphics[width=\linewidth]{p11_physical_delay.pdf}
\caption{El modelo físico de acumulación y su frontera de identificabilidad, sobre las crisis
NBER (log trimestre a trimestre, sin 2008/2020). \emph{Izquierda:} el ajuste ($R^2$ del modelo)
en función del retardo medio $\mu$ tiene un máximo nítido en $\mu\approx5$ trimestres ($\sim$1
año) ---el \emph{centro} del retardo se identifica. \emph{Centro:} ese mismo ajuste es plano
frente al ancho del núcleo (rango de $R^2\approx0{,}012$) ---el \emph{ancho} no se identifica.
\emph{Derecha:} dos núcleos Gamma con el mismo centro ($\mu=4$ trim) pero forma muy distinta
---casi puntual ($N=64$ etapas) y ancho ($N=2$)--- ajustan los datos por igual. El $R^2$ es el del
ajuste OLS del modelo al variar sus parámetros.}
\label{fig:model}
\end{figure}

El efecto es negativo, real y estructural. El coeficiente de sobreacumulación es \emph{negativo}:
una regresión de la ganancia sobre la
inversión madurada $m$ ---controlando por la inversión contemporánea y la ganancia rezagada--- da
$k\approx-0{,}40$, y el inverso bayesiano lo confirma con un IC95\% $[-0{,}63;-0{,}12]$ que no cruza
cero (Fig.~\ref{fig:bayes}). El tamaño del efecto se lee de este $k$ ---el coeficiente
\emph{condicional} del bayesiano, que descuenta la inversión contemporánea y la ganancia previa---,
no de la profundidad de la CCF: esa correlación marginal \emph{establece} el feedback y lo
\emph{ubica} en el rezago $-4$, pero no mide su magnitud estructural. Que sea negativo es el punto sustantivo: la inversión acumulada deprime
la ganancia posterior, capital que al acumularse erosiona su propia rentabilidad. El efecto es
estructural y no una firma de las crisis: el placebo compara la profundidad del valle en las
ventanas de crisis contra la de $\sim$2000 ventanas tomadas al azar fuera de ellas, y la joroba de
crisis \emph{no} resulta más profunda ($p\approx0{,}35$, una cola). Una firma de crisis habría dado
lo contrario ---un valle atípicamente profundo en las crisis, con $p$ chico---; que aparezca con la
misma fuerza fuera de ellas es la marca de un efecto siempre presente, no una falla de potencia del
test. Y no es sobreajuste: seis parámetros para 312 datos.

\begin{figure}[t]
\centering
\includegraphics[width=\linewidth]{p17_bayes_lag.pdf}
\caption{Inverso bayesiano del retardo (Turing.jl/NUTS, log-trim, 1948--2026). \emph{Izquierda:}
el posterior de $\mu$ (azul) se concentra cerca de $\sim$5 trimestres, mucho más angosto que el
prior (gris): el dato \emph{informa} el retardo. \emph{Derecha:} el posterior del coeficiente de
sobreacumulación $k$ queda a la izquierda de cero ---el efecto es significativo.}
\label{fig:bayes}
\end{figure}

\paragraph{Limitaciones.}
El feedback explica solo $\sim$3\% de la varianza del crecimiento de las ganancias, contra
$\sim$24\% del co-movimiento contemporáneo, y de esa debilidad se siguen dos límites. El primero es
que la \emph{dispersión} del retardo no se identifica: el perfil de verosimilitud en el ancho de la
distribución es plano y dos núcleos muy distintos ajustan igual (Fig.~\ref{fig:model}, centro y
der.), de modo que los datos no distinguen un retardo distribuido de uno puntual. El segundo es que
el mecanismo no se separa con estas dos series: un $\mathrm{VAR}(p\!\ge\!4)$ lineal reproduce el
mismo valle de la CCF y explica casi la misma varianza que nuestro modelo, así que el modelo de
acumulación es una reparametrización
interpretable de ese contenido lineal antes que un mecanismo separable de él ---en muestra. Fuera de
muestra, en cambio, los dos \emph{no} son equivalentes: el modelo estructural ata el centro del
núcleo $\mu$ a una cantidad medible aparte ---los tiempos de construcción y maduración sectoriales---,
una predicción que un VAR en forma reducida no puede hacer; es ahí, y no en el ajuste, donde el
mecanismo sería refutable. Y aun aceptando esa equivalencia en muestra, una
caída de la ganancia tras la inversión es compatible tanto con sobreacumulación (capacidad que un
mercado acotado no valida) como con un \emph{profit-squeeze} (la actividad eleva costos y comprime la
participación). Distinguir el canal exigiría datos de ventas, inventarios y utilización, ausentes en
este par de series.

\section{Discusión y conclusiones}

El cuadro sustantivo es consistente y honesto. \emph{Existe} un feedback negativo rezagado, \emph{consistente con} la sobreacumulación, con un
retardo de acumulación de alrededor de un año. Es estructural antes que específico de las crisis, ya
que el placebo encuentra el valle negativo con igual frecuencia fuera de ellas. Lo estimamos como un
núcleo de retardo distribuido, estimado con un inverso bayesiano y cotejado con la correlación cruzada. Lo que importa de ese feedback es su \emph{signo}, no su magnitud: un coeficiente consistentemente negativo es una presión de sobreacumulación persistente sobre la rentabilidad, débil en varianza pero estructural en su dirección. Dos
límites acotan la afirmación. Primero, el contenido económico ---un feedback de sobreacumulación con
las ganancias cediendo antes de las recesiones--- pertenece a la tradición previa (Astarita, 2012;
Tapia, 2023); nuestro aporte es el modelo y el análisis de identificabilidad, no el fenómeno.
Segundo, contamos con apenas nueve crisis de posguerra utilizables para estimar el retardo (de las
doce recesiones del NBER, descontando las extremas de 2008/2020 y los bordes de muestra). Con esa
muestra y una señal del $\sim$3\% de la varianza, los datos no pueden separar un retardo genuinamente
distribuido de uno puntual, ni el mecanismo de acumulación de un VAR lineal de memoria larga.

La clave de identificabilidad es esta: la señal débil ($\sim$3\%), con 312 trimestres, \emph{alcanza}
para establecer el \emph{efecto} ---un coeficiente negativo, real por el IC95\% bayesiano que no cruza
cero--- y su \emph{lag} de $\sim$1 año, pero \emph{no} para resolver lo fino: ni su \emph{dispersión}
(distribuido vs puntual) ni su \emph{separación de un VAR} se identifican. El resultado transferible
es \emph{dónde} cae esa frontera ---lo que la señal fija y lo que no---; cruzarla requiere datos
nuevos, no más método.

De cara al trabajo futuro, la señal es genuinamente escasa: EE.UU.\ es prácticamente la única serie
larga y confiable, y sus crisis están espaciadas $\sim$una por década, de modo que acumular más
historia rinde poco. Las vías realistas no pasan por más países sino por más resolución dentro del
mismo sistema: tiempos de construcción sectoriales medidos que predigan el núcleo macro ---un test
que un VAR en forma reducida no puede pasar--- o datos de mayor frecuencia. En el plano conceptual,
queda abierto cómo entra la \emph{resolución} de la crisis. Nuestro modelo lineal captura la
acumulación y su reflujo ---$m$ sube en el auge y decae al caer la inversión, el ida y vuelta que
subyace a la CCF oscilante---, pero representa ese reflujo como un desvanecimiento \emph{suave} de
$m$, no como la \emph{destrucción/desvalorización de capital} que descarga abruptamente la
sobreacumulación en la crisis. Esa descarga es no lineal y dependiente del estado; captarla
requeriría extender la cadena de acumulación con un mecanismo de crisis no lineal o de cambio de
régimen ---por ejemplo, una destrucción de capital por umbral o una saturación del efecto de
sobreacumulación---. Cualquier extensión de ese tipo debería juzgarse por comparación penalizada e
identificabilidad de sus parámetros, no por su ajuste en-muestra; anticipamos que, con la debilidad
de la señal ($\sim$3\%) y los nueve episodios, la mejora podría no resultar identificable, el mismo
límite que estructura este trabajo. Acoplar un ciclo generado por el retardo con una resolución de
crisis no lineal es, hasta donde sabemos, un problema abierto.

\section*{Apéndice: un inverso diferenciable (PINN) no mejora}
Probamos también un inverso con red sustituta (PINN), con las salvaguardas
estándar: red más chica que los datos, \emph{weight decay} y validación sobre una crisis retenida.
\emph{No} identifica el retardo: la red es flexible y ajusta los datos satisfaciendo la ODE para un
rango ancho de $\mu$, así que el costo queda casi plano en $\mu$ y la señal débil ($\sim$3\%) que
debería anclarlo se \emph{absorbe} en la red. El inverso bayesiano del cuerpo principal ---seis
parámetros directos, sin red--- deja que ese término débil informe $\mu$, y por eso \emph{sí} lo
identifica: con señal débil, la flexibilidad de la red destruye la identificabilidad que una
calibración probabilística directa preserva.

\section*{Agradecimientos}
El autor utilizó Claude (Anthropic) como asistente para redacción, figuras y código de los
experimentos. Todos los resultados, ecuaciones y conclusiones fueron verificados por el autor,
que asume la responsabilidad por el contenido.

\section*{Datos y software}
Los datos son series públicas de FRED: ganancias corporativas (\texttt{A053RC1Q027SBEA}),
inversión privada bruta (\texttt{GPDI}) y recesiones del NBER (\texttt{USREC}), con un snapshot
versionado para reproducir sin re-descargar. El código (Python~3.12 para el análisis empírico;
Julia con \texttt{Turing.jl} para el inverso bayesiano y \texttt{Lux.jl}/\texttt{Zygote.jl} para
la nota sobre el PINN) está en \url{https://github.com/JuaniTollo/-flight-to-quality}, con entornos
fijados (\texttt{uv.lock}, \texttt{Manifest.toml}). Series FRED de dominio público; librerías de
código abierto (MIT/BSD).

\section*{Referencias}
{\small
\setlength{\parindent}{0pt}
\setlength{\parskip}{4pt}
\hangindent=1.5em\hangafter=1 Astarita, R. (2012). Crisis, sobrecapacidad y coyuntura. \emph{rolandoastarita.blog}, 22 de marzo de 2012. \url{https://rolandoastarita.blog/2012/03/22/crisis-sobrecapacidad-y-coyuntura/}\par
\hangindent=1.5em\hangafter=1 Goldstein, J.P. (1999). Predator--prey model estimates of the cyclical profit squeeze. \emph{Metroeconomica} 50(2), 139--173.\par
\hangindent=1.5em\hangafter=1 Sapienza, F. (2026). \emph{Datos, Ecuaciones Diferenciales e Inteligencia Artificial} (apuntes de la materia, DM2026). \url{https://github.com/facusapienza21/DM2026-Curso}\par
\hangindent=1.5em\hangafter=1 Smith, H. (2011). \emph{An Introduction to Delay Differential Equations with Applications to the Life Sciences}. Springer. \url{https://doi.org/10.1007/978-1-4419-7646-8}\par
\hangindent=1.5em\hangafter=1 Strogatz, S.H. (2015). \emph{Nonlinear Dynamics and Chaos: With Applications to Physics, Biology, Chemistry, and Engineering}, 2.ª ed. Westview Press. ISBN 978-0-8133-4910-7.\par
\hangindent=1.5em\hangafter=1 Tapia, J.A. (2023). \emph{Six Crises of the World Economy: Globalization and Economic Turbulence from the 1970s to the COVID-19 Pandemic}. Palgrave Macmillan. \url{https://doi.org/10.1007/978-3-031-38735-7}\par
}

\end{document}
